\chapter*{Conclusione} %Se si cambia il Titolo cambiare anche la riga successiva così che appia corretto nell'conclusione
\addcontentsline{toc}{chapter}{Conclusione} %Per far apparire Introduzione nell'indice (Il nome deve rispecchiare quello del chapter)

L'integrazione dell'Open Source Intelligence con le risorse presenti nel Dark Web hanno offerto un approccio potente e innovativo alla protezione degli asset di un'organizzazione in un panorama di cyber threat sempre più complesso. Il caso studio proposto ha dimostrato come l'uso combinato di strumenti di monitoraggio e analisi, ambienti isolati e macchine dedicate, possa fornire un quadro completo per identificare, analizzare e mitigare le minacce.

Il Dark Web, con le sue caratteristiche di anonimato e difficile accessibilità, rappresenta una fonte unica di informazioni critiche per gli analisti di cybersecurity. Esso ha infatti permesso di accedere a dati e informazioni non indicizzati, di rivelare attività criminali, tentativi di vendita di credenziali compromesse, e altre minacce emergenti. In combinazione con le tecniche OSINT, che facendo uso di informazioni pubbliche, hanno permesso di ottenere un quadro chiaro di intelligence in grado di ampliare significativamente la capacità di individuazione delle minacce, comprendere le loro origini, i loro vettori di attacco, e quindi di reagire in modo tempestivo ed efficace\\

Il caso di studio ha dimostrato l'importanza dell'utilizzo del Dark Web come fonte di informazioni OSINT. Si è analizzato come un'organizzazione possa utilizzare determinate informazioni per migliorare le proprie abilità di difesa. L'implementazione di ambienti isolati, il monitoraggio continuo del Dark Web e l'adozione di contromisure appropriate sono stati elementi cruciali per garantire la sicurezza informatica e la protezione dei dati sensibili.

\'E stato analizzato come determinati strumenti siano fondamentali per il monitoraggio del Dark Web, per la creazione di regole di correlazione e creazione di "allarmi" di sicurezza, migliorando la capacità di risposta alle minacce. Inoltre, è stata evidenziata l'importanza di un'adeguata formazione per i dipendenti e di una continua revisione delle policy di sicurezza in previsione di possibili attacchi.

L'analisi combinata delle risorse è quindi fondamentale per la cybersecurity moderna. Mentre il Dark Web rappresenta un ambiente ricco di dati non convenzionali ma di grande valore, l'OSINT offre una base metodologica solida per raccogliere, analizzare e sfruttare questi dati. Tuttavia, operare in questi contesti comporta sfide significative, oltre a considerazioni etiche e legali che non devono essere trascurate.\\

In conclusione, l'integrazione tra questi due potenti strumenti rappresenta una strategia efficace e necessaria per le organizzazioni che desiderano rimanere resilienti contro le minacce. L'adozione di un approccio proattivo al monitoraggio dei cyber threat e alla risposta agli incidenti, unito a un costante aggiornamento delle tecniche e delle tecnologie, è essenziale per mitigare i rischi associati alle minacce zero-day e ad attacchi sempre più sofisticati. In un contesto digitale in continua evoluzione, la capacità di adattarsi rapidamente e di utilizzare tutte le risorse disponibili, sarà sempre più cruciale per la sicurezza e la protezione delle organizzazioni. 

